\documentclass[12pt,a4paper]{article}
\usepackage[UTF8]{ctex}          % 添加中文支持
\usepackage{amsmath,amssymb,amsfonts,geometry,enumitem}
\geometry{left=2.5cm,right=2.5cm,top=2.5cm,bottom=2.5cm}
\title{2024年新高考Ⅰ卷·数学 考场作答}
\author{考生：匿名\quad 准考证号：202406071234}
\date{2024年6月7日}
\begin{document}
\maketitle

\begin{center}
{\large \textbf{（本答案为模拟考生在考场上的真实作答情况，包含正确与错误作答，仅供参考）}}
\end{center}

\section*{一、选择题（1\~{}8题）}

\begin{enumerate}[label=\arabic*.]
\item \underline{A} \quad （这个会做，解不等式取交集就行）

\item \underline{C} \quad （复数的运算，花了点时间但算出来了）

\item \underline{D} \quad （向量垂直，数量积为0，算得$x=2$）

\item \underline{A} \quad （这题有点难，用两角和差公式展开，最后凑出来了）

\item \underline{B} \quad （圆柱和圆锥侧面积相等，列方程算体积，计算量有点大）

\item \underline{？} \quad （这题不太确定，蒙了个C，分段函数单调性没搞明白）

\item \underline{C} \quad （三角函数交点个数，画了草图数出来6个）

\item \underline{B} \quad （抽象函数递推，推了几项感觉选B，不知道对不对）
\end{enumerate}

\section*{二、多项选择题（9\~{}11题）}

\begin{enumerate}[label=\arabic*.]
\item \underline{BC} \quad （正态分布，$\mu=1.8$，$\sigma=0.1$，$P(X>2)=P(Z>2)\approx0.0228$，所以A错B对；$Y\sim N(2.1,0.01)$，$P(Y>2)=P(Z>-1)\approx0.8413$，所以C对D错）

\item \underline{ACD} \quad （求导分析单调性，A对；B代$x=0.5$试了一下发现不对；C和D用换元验证是对的——多选题不敢多选，就选了ACD）

\item \underline{？} \quad （这题完全没看懂，那个图也不知道是什么意思，随便蒙了BC）
\end{enumerate}

\section*{三、填空题（12\~{}14题）}

\begin{enumerate}[label=\arabic*.]
\item \underline{$\dfrac{3}{2}$} \quad （双曲线通径$|AB|=\dfrac{2b^2}{a}=10$，又$|F_1A|=13$，勾股定理列方程解得$e=\dfrac{3}{2}$——这题花了快10分钟）

\item \underline{$\ln 2$} \quad （切线斜率$k=2$，设切点，列方程解得$a=\ln 2$）

\item \underline{$\dfrac{7}{16}$} \quad （甲从8个数中选1个，乙从剩下7个中选1个，共$8\times7=56$种，甲大于乙的有$C_8^2=28$种，概率$\dfrac{28}{56}=\dfrac{1}{2}$——等等，我是不是算错了？应该是$\dfrac{7}{16}$？\dots 算了不改了，写$\dfrac{7}{16}$吧）
\end{enumerate}

\section*{四、解答题}

\begin{enumerate}[label=\arabic*.]
\item （13分）
\begin{enumerate}[label=(\arabic*)]
\item 由余弦定理：$\cos C=\dfrac{a^2+b^2-c^2}{2ab}=\dfrac{\sqrt{2}}{2}$，所以$C=45^\circ$。

又$\sin C=\sqrt{2}\cos B$，$\dfrac{\sqrt{2}}{2}=\sqrt{2}\cos B$，$\cos B=\dfrac{1}{2}$，$B=60^\circ$。

\item $A=75^\circ$。面积$S=\dfrac{1}{2}ac\sin B=3+\sqrt{3}$，$\dfrac{\sqrt{3}}{4}ac=3+\sqrt{3}$，$ac=4\sqrt{3}+4$。

由正弦定理$\dfrac{a}{\sin75^\circ}=\dfrac{c}{\sin45^\circ}$，$a=\dfrac{c\sin75^\circ}{\sin45^\circ}$。

代入解得$c=2\sqrt{2}$。（这题做得还算顺，希望过程分能拿满）
\end{enumerate}

\item （15分）
\begin{enumerate}[label=(\arabic*)]
\item $b^2=9$，代入$P$点得$a^2=12$，$e=\dfrac{1}{2}$。

\item 设直线$l:y-\dfrac{3}{2}=k(x-3)$。点$A$到$l$的距离$d=\dfrac{|3k-\frac{9}{2}|}{\sqrt{k^2+1}}$，$|AP|=\dfrac{3\sqrt{5}}{2}$。

$S=\dfrac{1}{2}\cdot\dfrac{3\sqrt{5}}{2}\cdot d=9$，解得$k=\dfrac{3}{2}$或$k=-\dfrac{3}{2}$。

所以$l:y=\dfrac{3}{2}x-3$或$y=-\dfrac{3}{2}x+6$。——（这题计算量真大，算了两遍才确定答案）
\end{enumerate}

\item （15分）
\begin{enumerate}[label=(\arabic*)]
\item 因为$PA\perp$底面$ABCD$，所以$PA\perp AD$。又$AD\perp PB$，且$PA\cap PB=P$，所以$AD\perp$平面$PAB$，故$AD\perp AB$。又$AD\subset$平面$ABCD$，\dots （这题写了一半发现证不下去了，先跳过去做后面的）

（留了大片空白，回头再补——最后也没补上，立体几何一直是我的弱项）

\item 建系做了半天，算出来$AD=3$，不知道对不对。
\end{enumerate}

\item （17分）
\begin{enumerate}[label=(\arabic*)]
\item $f'(x)=\dfrac{2}{x(2-x)}+a\geqslant0$，$a\geqslant-\dfrac{2}{x(2-x)}$。$x(2-x)\leqslant1$，所以$a\geqslant-2$。最小值为$-2$。

\item 中心对称图形？$f(x)+f(2-x)=2a$，所以关于$(1,a)$中心对称。（这题意外的简单）

\item 不会做，只写了$b$的范围大概是$\left[-\dfrac{2}{3},+\infty\right)$，过程没来得及写完。
\end{enumerate}

\item （17分）
\begin{enumerate}[label=(\arabic*)]
\item $(i,j)=(1,2),(1,6),(5,6)$。——（枚举了所有情况，应该没错）

\item 不会证明，只写了一句话“由等差数列性质可构造分组”就放弃了。

\item 更不会，直接空着了。
\end{enumerate}

\vspace{1cm}

\begin{flushright}
\textbf{【考场状态备注】}\\
整体感觉：比平时模考难很多，选填花了太多时间，\\
解答题第15题立体几何卡住了，第19题完全不会。\\
估计分数：90\~{}100分左右。\\
——来自一位真实考生的考场回忆
\end{flushright}

\end{enumerate}
\end{document}



